% Werke X/1 (1917), native170–171 upper, printed163–164 upper.
\begingroup
\setlength{\headheight}{25pt}
\SourceLeaf{170}\fancyhead{}\fancyhead[R]{\small163}\fancyhead[C]{\scriptsize\shortstack{ÄLTESTE UNTERSUCHUNGEN ÜBER\\ LEMNISKATISCHE FUNKTIONEN.\quad III.}}
Hieraus aber folgt eine zweite Wurzel
\[=16\frac{(1-t)^2t}{(1+t)^2}\]
{\small\[\renewcommand{\arraystretch}{1.4}\begin{array}{rrl}
l16={}&1{,}2041200&\\
lt={}&0{,}8655836-2&\begin{aligned}&[=\log_{10}0{,}073381\ldots\\&=\log_{10}(\sin\text{ lemn }36^0)^4]\end{aligned}\\
l(1-t)^2={}&0{,}9338025-1&\\
-l(1+t)^2={}&0{,}8769844-1&\\\cline{2-2}
&0{,}8804905-1&\begin{aligned}&[=\log_{10}]\ 0{,}7594348\\&[=\log_{10}(\sin\text{ lemn }72^0)^4].\end{aligned}
\end{array}\]}
\noindent Also
\[l\sin72^0=0{,}9701226\]
\noindent und
\[\sin72^0=0{,}9335179.\]
\NumberedItem{6}
\HistoricCitation{[Ein Zettel Fh Nr. 1, Kapsel 50.]}
\HistoricTitle{Rechnungen zur Lemniscata gehörig.}
{\small\[\begin{gathered}
(\sin0{,}4)^4+(\sin0{,}8)^4=12\surd5-26\ [^{*)}]\\
=2(0{,}4164078649\quad9873817845\quad5042012387\quad65741)
\end{gathered}\]}
{\small\[\begin{gathered}
\bigl(\tfrac12(\sin0{,}8)^4-\tfrac12(\sin0{,}4)^4\bigr)^2=340-152\surd5\\
=0{,}1176674200\quad3196614580\quad5602352846\quad01228.
\end{gathered}\]}
\noindent Daraus Radix
\[\begin{array}{rrrr}
0{,}3430268503&0761971797&7310507555&85731\\
&&-0463415326&98758\\\cline{3-4}
&&6847092228&86973
\end{array}\]
\noindent Also
\[(\sin0{,}4)^4=0{,}0733810146\quad9111846047\quad7731504831\quad80010\]
\noindent Daraus Radix
\[(\sin0{,}4)^2=0{,}2708893033\quad8999814497\quad30710\]
\renewcommand{\thefootnote}{[*)}\footnotetext{\(\sin0{,}4\) und \(\sin0{,}8\) bedeuten \(\sin\frac{4\varpi}{10}\) und \(\sin\frac{8\varpi}{10}\), d. h. also \(\sin\text{ lemn }\frac{2\varpi}{5}\) und \(\sin\text{ lemn }\frac{4\varpi}{5}\).]\par\medskip\noindent\hfill{\normalsize21\textsuperscript{*}}}
\clearpage
\SourceLeaf{171}\ArithmeticHead{164}{L}{ANALYSIS.\quad NACHLASS.}
\[(\sin0{,}8)^4=0{,}7594347153\quad0635789643\quad2352519994\quad51472\]
\noindent Daraus Radix
\[(\sin0{,}8)^2=0{,}8714555153\quad9155336074\quad646029.\]
\NumberedItem{7}
\HistoricCitation{[In \textit{Sammlung von Rechnungen usw.} Fh Nr. 2, Kapsel 50.]}
\HistoricTitle{Berechnungen die Lemniscata betreffend.}
\noindent\((\sin0{,}4)^4\) u. \((\sin0{,}8)^4\) sind Wurzeln der Gleichung[en]
\[\begin{aligned}
p^4+q^4&=\surd720-26\\
(pp+qq)^2&=14\surd5-30\\
(pp-qq)^2&=10\surd5-22\\
pp+qq&=5^{\frac14}(3-5^{\frac12})
\end{aligned}\]
\endgroup
